\documentclass[11pt,a4paper]{article}

\usepackage[margin=1in]{geometry}
\usepackage[T1]{fontenc}
\usepackage[utf8]{inputenc}
\usepackage{lmodern}
\usepackage{microtype}
\usepackage{amsmath,amssymb,mathtools}
\usepackage{booktabs,tabularx,array}
\usepackage{enumitem}
\usepackage{hyperref}
\usepackage{needspace}

\hypersetup{
  hidelinks
}

\setlength{\parindent}{0pt}
\setlength{\parskip}{0.55em}
\setlist[itemize]{leftmargin=1.4em}
\newcolumntype{Y}{>{\raggedright\arraybackslash}X}
\setlength{\emergencystretch}{3em}
\allowdisplaybreaks
\newcommand{\R}{\mathbb{R}}

\begin{document}

\begin{center}
{\Large\bfseries A Grad-TTS-Style Diffusion Acoustic Model}\\[0.5em]
{\large Architecture, Probabilistic Formulation, Losses, and Context in Modern TTS}
\end{center}

\begin{abstract}
This note describes a compact Grad-TTS-style acoustic model for text-to-mel synthesis. The model is non-autoregressive: text is encoded in parallel, durations determine frame expansion, a conditional prior predicts a coarse mel representation, and a diffusion decoder refines this prior into the final mel-spectrogram. Relative to the original Grad-TTS formulation, the present model differs in three main ways: alignment supervision is provided externally rather than learned by monotonic alignment search, the mel prior is produced by a separate prior network, and diffusion training includes an auxiliary clean reconstruction term.
\end{abstract}

\section{Overview}

The original Grad-TTS model of Popov et al.~\cite{popov2021gradtts} is a non-autoregressive acoustic model: an encoder maps text to a frame-aligned Gaussian mean, a duration model determines token expansion, and a score-based decoder gradually denoises a noisy sample into a mel-spectrogram. The present model follows the same general idea but is not a literal reproduction of the paper.

The main conceptual point is that this model does \emph{not} learn monotonic alignments internally with Monotonic Alignment Search (MAS), unlike Grad-TTS and Glow-TTS~\cite{kim2020glowtts,popov2021gradtts}. Instead, training uses externally prepared token durations, so alignment learning lies outside the acoustic model itself.

The model is also an \emph{acoustic} model rather than a full waveform synthesizer. Its output is a normalized mel-spectrogram, so a separate waveform reconstruction stage is still required.

Within the broader TTS landscape, one can distinguish autoregressive models such as Tacotron~2~\cite{shen2018naturaltts}, non-autoregressive duration-based models such as FastSpeech and FastSpeech~2~\cite{ren2019fastspeech,ren2020fastspeech2}, flow-based systems such as Glow-TTS~\cite{kim2020glowtts}, diffusion-based models such as Grad-TTS and NaturalSpeech~2~\cite{popov2021gradtts,shen2023naturalspeech2}, end-to-end variational-adversarial systems such as VITS~\cite{kim2021vits}, and codec language model approaches such as VALL-E~\cite{wang2023valle}. The present model is conceptually closest to the duration-based non-autoregressive line, but it replaces deterministic acoustic decoding with a diffusion decoder.

\section{Architecture}

\subsection{Concrete Specification}

The instantiated model uses the structural hyperparameters shown in Table~\ref{tab:config}.

\begin{table}[h]
\centering
\caption{Main architectural specification of the present model.}
\label{tab:config}
\begin{tabularx}{\linewidth}{@{}lY@{}}
\toprule
Parameter & Value \\
\midrule
Mel bins $M$ & $80$ \\
Encoder hidden size $H$ & $512$ \\
Encoder prenet layers & $3$ conv layers, kernel size $5$ \\
Encoder blocks & $6$ \\
Attention heads & $8$ \\
Encoder FFN expansion & $4\times$ (so $512 \rightarrow 2048 \rightarrow 512$) \\
Duration predictor & $2$ conv layers with width $256$, kernel size $3$ \\
Prior predictor & $3$ residual conv blocks, kernel size $5$ \\
Diffusion U-Net base width & $64$ \\
Diffusion channel multipliers & $(1,2,4)$ \\
Diffusion steps at inference & $80$ \\
Noise schedule & $\beta_{\min}=0.05$, $\beta_{\max}=20.0$, $t\in[0,1]$ \\
Training crop & $192$ mel frames \\
Default sampling mode & Deterministic reverse diffusion \\
Temperature parameterization & Inverse temperature \\
\bottomrule
\end{tabularx}
\end{table}

\subsection{Layer Architecture}

\begin{itemize}
\item \textbf{Text encoder.} Text tokens are embedded into a $512$-dimensional space. The embedding is processed by a three-layer convolutional prenet, followed by six Transformer-like encoder blocks. Each block consists of relative multi-head self-attention and a convolutional feed-forward subnetwork, both with residual connections and normalization.

\item \textbf{Duration predictor.} A two-layer temporal convolutional network maps token-level encoder states to one scalar per token. This scalar represents the logarithm of token duration in mel frames.

\item \textbf{Length regulator.} Token-level hidden states are repeated according to integer durations, producing frame-level conditioning. This converts the model from text length $L$ to acoustic length $F$.

\item \textbf{Prior network.} The frame-level conditioning is transformed into a coarse mel representation $\mu \in \R^{80\times F}$ by a residual convolutional network. This $\mu$ acts as the mean of the text-conditioned Gaussian prior.

\item \textbf{Diffusion decoder.} The decoder is a U-Net score model. Its input is a two-channel tensor formed by stacking the current noisy spectrogram state $x_t$ with the prior mean $\mu$. Time is encoded sinusoidally, and frame-level conditioning is injected at several scales of the U-Net.

\item \textbf{U-Net channel progression.} The down path uses widths $64 \rightarrow 128 \rightarrow 256$, the bottleneck remains at width $256$, and the up path mirrors the hierarchy with skip connections. Each scale uses residual blocks and lightweight attention.
\end{itemize}

\subsection{Forward Pipeline}

The overall acoustic pipeline can be written as
\[
\begin{aligned}
x_{1:L}
&\longrightarrow
h_{1:L}\in\R^{L\times 512}
\longrightarrow
\hat{\ell}_{1:L}
\longrightarrow
c_{1:F}\in\R^{F\times 512}
\\
&\longrightarrow
\mu\in\R^{80\times F}
\longrightarrow
\hat{x}_0\in\R^{80\times F}.
\end{aligned}
\]

At training time, the model uses externally supplied durations $d_{1:L}$ to build frame-level conditioning. At inference time, durations are predicted, rounded to integers, and then used for length regulation before diffusion-based mel generation.

\section{Probabilistic Formulation and Losses}

\subsection{Notation}

\Needspace{10\baselineskip}
\begin{tabularx}{\linewidth}{@{}lY@{}}
\toprule
Symbol & Meaning \\
\midrule
$x_{1:L}$ & text token sequence of length $L$ \\
$h_{1:L}$ & encoder hidden states \\
$d_{1:L}$ & token durations measured in mel frames \\
$c_{1:F}$ & frame-level conditioning after length regulation \\
$x_0 \in \R^{M\times F}$ & target normalized mel-spectrogram \\
$\mu \in \R^{M\times F}$ & prior mean predicted from the conditioning sequence \\
$M$ & number of mel bins, here $M=80$ \\
$F$ & number of acoustic frames \\
$L$ & number of text tokens \\
$m_f \in \{0,1\}$ & frame-validity mask at frame $f$ \\
$N$ & number of valid mel entries, $N=M\sum_{f=1}^{F} m_f$ \\
\bottomrule
\end{tabularx}

In the present model, the time variable is normalized to $t\in[0,1]$ and the linear noise schedule is
\[
\beta(t) = \beta_{\min} + (\beta_{\max} - \beta_{\min})t,
\qquad
\bar{\beta}(t) = \int_0^t \beta(s)\,ds
= \beta_{\min} t + \frac{1}{2}(\beta_{\max}-\beta_{\min})t^2.
\]
Define
\[
\alpha_t = \exp\!\left(-\frac{1}{2}\bar{\beta}(t)\right),
\qquad
\lambda_t = 1 - \exp(-\bar{\beta}(t)) = 1-\alpha_t^2.
\]
Here, $\beta_{\min}$ and $\beta_{\max}$ are the minimum and maximum noise rates, $\bar{\beta}(t)$ is the cumulative noise level, $\alpha_t$ is the clean-signal shrinkage factor, and $\lambda_t$ is the variance of the corrupted sample at time $t$.

\subsection{Prior Model}

The present model has an explicit prior network:
\[
\mu = P_{\phi}(c).
\]
The prior loss is
\[
\mathcal{L}_{\text{prior}}
= \frac{1}{N}
\sum_{f=1}^{F}\sum_{m=1}^{M}
\frac{1}{2}\left((x_{0,mf}-\mu_{mf})^2 + \log(2\pi)\right)m_f,
\]
where $x_{0,mf}$ is the target mel value at mel bin $m$ and frame $f$, $\mu_{mf}$ is the corresponding prior mean, and $m_f$ masks out padded frames. The term $\log(2\pi)$ is the constant part of the negative log-likelihood of a unit-variance Gaussian; it appears because the prior is written as a full Gaussian likelihood rather than only as a squared-error term.

This is exactly the negative log-likelihood of an isotropic Gaussian
\[
p_{\phi}(x_0 \mid c) = \mathcal{N}(x_0;\mu_{\phi}(c), I),
\]
computed framewise in mel space.

\subsection{Forward Diffusion}

The forward corruption process is
\[
q(x_t \mid x_0,\mu)
=
\mathcal{N}\!\left(
x_t;\;
\alpha_t x_0 + (1-\alpha_t)\mu,\;
\lambda_t I
\right).
\]
Equivalently, with $\varepsilon\sim\mathcal{N}(0,I)$,
\[
x_t = \alpha_t x_0 + (1-\alpha_t)\mu + \sqrt{\lambda_t}\,\varepsilon.
\]
Here, $x_t$ is the noisy spectrogram state at time $t$, and $\varepsilon$ is standard Gaussian noise.

This is the key Grad-TTS idea: diffusion is not centered at zero, but at the prior mean $\mu$. Relative to standard DDPMs~\cite{ho2020ddpm}, the model starts from a text-conditioned Gaussian rather than from an unconditional standard normal.

\subsection{Score Network Objective}

The estimator $s_{\theta}(x_t,\mu,c,t)$ is trained as a score model. For the Gaussian above, the true conditional score is
\[
\nabla_{x_t} \log q(x_t \mid x_0,\mu)
=
-\frac{x_t - \alpha_t x_0 - (1-\alpha_t)\mu}{\lambda_t}
= -\frac{\varepsilon}{\sqrt{\lambda_t}}.
\]

The score loss is
\[
\mathcal{L}_{\text{score}}
= \frac{1}{N}
\left\|
\sqrt{\lambda_t}\, s_{\theta}(x_t,\mu,c,t) + \varepsilon
\right\|_2^2.
\]
Here, $s_{\theta}(x_t,\mu,c,t)$ is the learned score field and $N$ is the normalization by the number of valid mel entries.
This matches the score-matching form used by Grad-TTS~\cite{popov2021gradtts}, written in the normalization adopted by the present model.

\subsection{Auxiliary Clean Reconstruction Loss}

Unlike the original paper, the present model adds an auxiliary clean-signal term. From the score estimate, it reconstructs
\[
\hat{x}_0
=
\frac{x_t - (1-\alpha_t)\mu + \lambda_t s_{\theta}(x_t,\mu,c,t)}{\alpha_t}.
\]
Then it computes an L1 loss
\[
\mathcal{L}_{x_0}
=
\frac{1}{N}\left\|\hat{x}_0 - x_0\right\|_1.
\]
Here, $\hat{x}_0$ is the clean spectrogram reconstructed from the noisy sample and the predicted score.

The diffusion loss is therefore
\[
\mathcal{L}_{\text{diff}}
=
\mathcal{L}_{\text{score}} + \gamma \mathcal{L}_{x_0},
\]
where $\gamma$ is the weight of the auxiliary reconstruction term. In the present configuration, $\gamma=0.05$.

\subsection{Duration Loss}

The duration head predicts a scalar log-duration per token. The target is
\[
\ell_i = \log(1 + d_i),
\]
where $d_i$ is the ground-truth duration of token $i$ in frames and $\ell_i$ is its log-transformed target. The logarithm compresses the wide dynamic range of durations and makes regression more stable.

The loss is masked MSE:
\[
\mathcal{L}_{\text{dur}}
=
\frac{1}{\sum_i u_i}
\sum_{i=1}^{L} u_i \big(\hat{\ell}_i - \log(1+d_i)\big)^2,
\]
where $i$ indexes text tokens, $L$ is the number of tokens, $\hat{\ell}_i$ is the predicted log-duration, and $u_i \in \{0,1\}$ is a token-validity mask that ignores padding positions.

This resembles the duration regression objective used in FastSpeech and Grad-TTS~\cite{ren2019fastspeech,popov2021gradtts}, but here the supervision comes from externally supplied durations rather than MAS estimated inside the model.

\subsection{Total Training Objective}

The total objective is
\[
\mathcal{L}_{\text{total}}
=
\mathcal{L}_{\text{diff}}
+
w_{\text{prior}}\mathcal{L}_{\text{prior}}
+
w_{\text{dur}}\mathcal{L}_{\text{dur}},
\]
with
\[
w_{\text{prior}} = 1.0,
\qquad
w_{\text{dur}} = 0.05.
\]
Here, $w_{\text{prior}}$ and $w_{\text{dur}}$ control the relative contribution of the prior and duration objectives.

An additional prior L1 quantity may be monitored diagnostically, but it is not part of the optimized objective.

\section{Generation and Implied Distributions}

\subsection{Text to Conditioning}

At inference, text is encoded as
\[
h_{1:L} = E_{\psi}(x_{1:L}).
\]
The duration head predicts log-durations $\hat{\ell}_{1:L}$, which are converted to frame counts by
\[
\hat{d}_i
=
\operatorname{clip}
\left(
\operatorname{round}\!\left(
\max(\exp(\hat{\ell}_i)-1,0)\cdot s
\right),
0,d_{\max}
\right),
\]
where $\hat{\ell}_i$ is the predicted log-duration, $\exp(\hat{\ell}_i)-1$ inverts the training transform $\log(1+d_i)$, $\max(\cdot,0)$ prevents negative durations, $s$ is a length-scale factor controlling speaking rate, $\operatorname{round}(\cdot)$ converts the continuous estimate into an integer number of frames, and $\operatorname{clip}(\cdot,0,d_{\max})$ prevents pathological overly long durations. In brief, this formula turns a real-valued log-duration prediction into a valid, bounded frame count suitable for length regulation. After this, length regulation produces
\[
c = \operatorname{LR}(h,\hat{d}),
\qquad
\mu = P_{\phi}(c).
\]

\subsection{Terminal Distribution}

The initial noisy latent is sampled with inverse-temperature scaling:
\[
z_1 \sim p(z_1 \mid \mu)
=
\mathcal{N}\!\left(\mu,\tau^{-2}I\right),
\]
equivalently,
\[
z_1 = \mu + \frac{\varepsilon}{\tau}, \qquad \varepsilon\sim\mathcal{N}(0,I).
\]
Here, $\tau$ is the temperature parameter. Larger $\tau$ means weaker initial noise under this inverse-temperature parameterization.

This detail is easy to miss: with inverse temperature mode, a \emph{larger} temperature means \emph{less} initial noise. That is the opposite of the more common $z=\mu+\tau\varepsilon$ parameterization.

\subsection{Reverse Diffusion Update}

The reverse process is solved numerically with $N$ Euler steps. In the present configuration, $N=80$.

For $\Delta t = 1/N$ and midpoint time
\[
t_k = 1 - \left(k+\frac{1}{2}\right)\Delta t,
\]
the deterministic update used by default is
\[
x_{t_k-\Delta t}
=
x_{t_k}
-
\left(
\frac{1}{2}(\mu - x_{t_k}) - s_{\theta}(x_{t_k},\mu,c,t_k)
\right)\beta(t_k)\Delta t.
\]
Here, $k$ indexes reverse steps, $\Delta t = 1/N$ is the step size, and $\beta(t_k)$ is the instantaneous noise rate at step $k$.

If stochastic sampling is enabled, the model adds a Gaussian perturbation with variance $\beta(t_k)\Delta t$ at each step. The final sample after the last step is interpreted as the predicted normalized mel-spectrogram $\hat{x}_0$.

\subsection{End-to-End Generative View}

Conditioned on text, the implemented model induces the following implicit generative chain:
\[
\begin{aligned}
x_{1:L}
&\longrightarrow
h_{1:L}
\longrightarrow
\hat{d}_{1:L}
\longrightarrow
c_{1:F}
\longrightarrow
\mu
\\
&\longrightarrow
z_1 \sim \mathcal{N}(\mu,\tau^{-2}I)
\longrightarrow
\hat{x}_0 = \operatorname{RevDiff}_{\theta}(z_1,\mu,c).
\end{aligned}
\]

Therefore a compact probabilistic summary is
\[
\begin{aligned}
p_{\theta,\phi,\psi}(\hat{x}_0 \mid x)
&=
\int
p_{\theta}(\hat{x}_0 \mid z_1,\mu,c)\,
p(z_1 \mid \mu)
\\
&\qquad \cdot\,
\delta\!\big(\mu-P_{\phi}(c)\big)\,
\delta\!\big(c-\operatorname{LR}(E_{\psi}(x),\hat{d})\big)\,
dz_1.
\end{aligned}
\]
The two delta terms reflect the fact that the encoder, duration expansion, and prior predictor are deterministic once text and predicted durations are fixed.

\section{Relation to the Original Grad-TTS Paper}

The main differences from the original Grad-TTS paper are the following.

\begin{itemize}
\item \textbf{Alignment learning.} Original Grad-TTS learns monotonic alignments with MAS~\cite{popov2021gradtts,kim2020glowtts}. The present model assumes durations are available during training, so alignment estimation is external.

\item \textbf{Prior parameterization.} In the original formulation, alignment and prior construction are tightly tied to the encoder representation. Here, a separate prior network maps frame-level conditioning to the mel-space mean $\mu$.

\item \textbf{Diffusion objective.} The core score-matching idea remains unchanged, but the present model adds an auxiliary clean reconstruction term $\mathcal{L}_{x_0}$.

\item \textbf{Conditioning strategy.} The diffusion U-Net is conditioned not only on $\mu$ and time, but also on frame-level hidden features injected across scales.

\item \textbf{Model scope.} The model remains an acoustic model that predicts mel-spectrograms, so waveform synthesis is delegated to a separate stage.
\end{itemize}

Accordingly, the most precise conceptual description is that this is a compact \emph{Grad-TTS variant} with externally supervised durations, an explicit mel prior network, and a diffusion decoder trained with both score loss and auxiliary clean reconstruction.

\section{Results}

This section summarizes the saved training history and the held-out inference metrics.

The model was trained for $80$ epochs. Training loss decreases from $2.1426$ at epoch $1$ to $1.1034$ at epoch $80$, showing stable optimization. The best validation checkpoint is obtained earlier, at epoch $41$, with
\[
\mathcal{L}_{\text{val}} = 1.2219,
\qquad
\mathcal{L}_{\text{val,diff}} = 0.0871,
\qquad
\mathcal{L}_{\text{val,dur}} = 0.0947,
\qquad
\mathcal{L}_{\text{val,prior}} = 1.1300,
\]
while the best training loss appears later at epoch $77$. The plots therefore suggest mild overfitting in the second half of training.

Held-out inference was evaluated on $12$ utterances. The mean metrics are
\[
\text{mel-MAE} = 3.6653,
\qquad
\text{mel-RMSE} = 4.5824,
\qquad
\text{duration-MAE} = 1.7875,
\qquad
\text{RTF} = 0.4308.
\]
The average predicted length is $622.8$ frames versus $664.3$ reference frames, so the model tends to generate slightly shorter utterances. The best sample in the recorded subset is \texttt{LJ003-0269} with mel-MAE $2.9671$, while the weakest is \texttt{LJ003-0096} with mel-MAE $4.6863$. The plots therefore indicate that the model learns the general acoustic structure, but quality still varies noticeably across sentences.

\section{Discussion}

The model is able to learn a stable text-to-mel mapping, but the held-out results show that it still underestimates duration and does not generalize as well as the training loss alone might suggest. The dependence on external alignments is also a clear limitation, since duration quality is tied directly to the preprocessing pipeline.

Future work should focus on stronger evaluation and stronger generation. In particular, it would be useful to add perceptual and ASR-based metrics, improve duration modeling, and test larger or better-conditioned diffusion backbones.

\section{Contributions}

\textbf{Sofin Mikhail.} Sofin Mikhail prepared the dataset pipeline, implemented data loading, and developed the encoder side of the model. He also contributed to early concept testing and analysis of intermediate results, helping establish the practical training setup.

\textbf{Vlad Strelkov.} Vlad Strelkov implemented the diffusion component and carried out iterative improvements of the overall model behavior. He also performed repeated personal testing and correction of the generated outputs, which guided the later refinement of the system.

\begin{thebibliography}{99}

\bibitem{ho2020ddpm}
Jonathan Ho, Ajay Jain, and Pieter Abbeel.
\newblock \emph{Denoising Diffusion Probabilistic Models}.
\newblock arXiv preprint arXiv:2006.11239, 2020.
\newblock \url{https://arxiv.org/abs/2006.11239}

\bibitem{kim2020glowtts}
Jaehyeon Kim, Sungwon Kim, Jungil Kong, and Sungroh Yoon.
\newblock \emph{Glow-TTS: A Generative Flow for Text-to-Speech via Monotonic Alignment Search}.
\newblock In \emph{NeurIPS}, 2020.
\newblock \url{https://arxiv.org/abs/2005.11129}

\bibitem{kim2021vits}
Jaehyeon Kim, Jungil Kong, and Juhee Son.
\newblock \emph{Conditional Variational Autoencoder with Adversarial Learning for End-to-End Text-to-Speech}.
\newblock In \emph{ICML}, 2021.
\newblock \url{https://arxiv.org/abs/2106.06103}

\bibitem{popov2021gradtts}
Vadim Popov, Ivan Vovk, Vladimir Gogoryan, Tasnima Sadekova, and Mikhail Kudinov.
\newblock \emph{Grad-TTS: A Diffusion Probabilistic Model for Text-to-Speech}.
\newblock In \emph{Proceedings of the 38th International Conference on Machine Learning (PMLR 139)}, 2021.
\newblock \url{https://arxiv.org/abs/2105.06337}

\bibitem{ren2019fastspeech}
Yi Ren, Yangjun Ruan, Xu Tan, Tao Qin, Sheng Zhao, Zhou Zhao, and Tie-Yan Liu.
\newblock \emph{FastSpeech: Fast, Robust and Controllable Text to Speech}.
\newblock In \emph{NeurIPS}, 2019.
\newblock \url{https://arxiv.org/abs/1905.09263}

\bibitem{ren2020fastspeech2}
Yi Ren, Chenxu Hu, Xu Tan, Tao Qin, Sheng Zhao, Zhou Zhao, and Tie-Yan Liu.
\newblock \emph{FastSpeech 2: Fast and High-Quality End-to-End Text to Speech}.
\newblock In \emph{ICLR}, 2021.
\newblock \url{https://arxiv.org/abs/2006.04558}

\bibitem{shen2018naturaltts}
Jonathan Shen, Ruoming Pang, Ron J. Weiss, Mike Schuster, Navdeep Jaitly, Zongheng Yang, Zhifeng Chen, Yu Zhang, Yuxuan Wang, RJ Skerry-Ryan, Rif A. Saurous, Yannis Agiomyrgiannakis, and Yonghui Wu.
\newblock \emph{Natural TTS Synthesis by Conditioning WaveNet on Mel Spectrogram Predictions}.
\newblock In \emph{ICASSP}, 2018.
\newblock \url{https://arxiv.org/abs/1712.05884}

\bibitem{shen2023naturalspeech2}
Kai Shen, Zeqian Ju, Xu Tan, Yanqing Liu, Yichong Leng, Lei He, Tao Qin, Sheng Zhao, and Jiang Bian.
\newblock \emph{NaturalSpeech 2: Latent Diffusion Models are Natural and Zero-Shot Speech and Singing Synthesizers}.
\newblock arXiv preprint arXiv:2304.09116, 2023.
\newblock \url{https://arxiv.org/abs/2304.09116}

\bibitem{song2021sde}
Yang Song, Jascha Sohl-Dickstein, Diederik P. Kingma, Abhishek Kumar, Stefano Ermon, and Ben Poole.
\newblock \emph{Score-Based Generative Modeling through Stochastic Differential Equations}.
\newblock In \emph{ICLR}, 2021.
\newblock \url{https://arxiv.org/abs/2011.13456}

\bibitem{wang2023valle}
Chengyi Wang, Sanyuan Chen, Yu Wu, Ziqiang Zhang, Long Zhou, Shujie Liu, Zhuo Chen, Yanqing Liu, Huaming Wang, Jinyu Li, Lei He, Sheng Zhao, and Furu Wei.
\newblock \emph{Neural Codec Language Models are Zero-Shot Text to Speech Synthesizers}.
\newblock arXiv preprint arXiv:2301.02111, 2023.
\newblock \url{https://arxiv.org/abs/2301.02111}

\end{thebibliography}

\end{document}
